\documentclass[UTF8]{ctexart}
\usepackage{geometry}
\usepackage{hyperref}
\usepackage{listings}
\usepackage{xcolor}

\geometry{a4paper, left=2.5cm, right=2.5cm, top=2.5cm, bottom=2.5cm}

\title{Agent开发周报 260422}
\author{刘德辰}
\date{2026年4月22日}

\begin{document}

\maketitle

\section{报告管线继续统一到 MCP 来源：线路、单位、车辆、驾驶员的来源结构和指标树对齐}

\textbf{背景}：上周已经把驾驶员、单位、线路、事故等报告推进到接近车辆报告的完成状态，本周重点从「能生成」转向「来源可信、结构一致、可回归」。报告如果只在正文里看起来完整，但 \texttt{report\_source}、指标树、建议来源和趋势来源不一致，后续前端展示、追溯、手测和线上排障都会出现口径分裂。

\textbf{方案说明}：本周围绕 MCP report source 做了多轮集中完善：统一结构化 lookup 使用 MCP 来源，补齐线路/单位/驾驶员/车辆的建议类 MCP 接入，抽出画像指标树处理能力，约束摘要指标必须来自叶子节点，并给报告来源增加契约脚本，防止后续适配层再次泄漏原始 MCP 包络或混用非叶子指标。

\textbf{功能说明}：
\begin{itemize}
\item 报告内容的一致性明显提升：不同对象（车辆、驾驶员、单位、线路）的报告结构更统一，阅读和对比更顺畅。
\item 建议内容更完整：报告中的改进建议覆盖面更全，空白项和「无内容占位」情况减少。
\item 趋势解读更稳定：同一类指标在不同报告中的解释更一致，用户不容易出现「同指标不同说法」的困惑。
\item 重点风险更靠前：高风险项展示顺序更符合直觉，用户打开报告后更容易先看到关键问题。
\item 报告可信度更高：来源口径做了统一梳理，减少「看起来像结论但依据不清晰」的体验。
\end{itemize}

\textbf{相关提交（日期）}：
\begin{itemize}
\item \texttt{3ea206f}（2026-04-21）add report source contract
\item \texttt{9a6910a}（2026-04-21）align report quota tree sources
\item \texttt{a6a93a8}（2026-04-19）wire report suggestion mcp sources
\item \texttt{b147f4f}（2026-04-18）unify report lookup on mcp sources
\item \texttt{be26c9a}（2026-04-20）Fix report dashboard risk ordering
\end{itemize}

\section{Router 与召回继续收紧：报告、专家、咨询边界更清楚，失败后能回到正确路径}

\textbf{背景}：多轮会话里最容易出问题的不是单次生成，而是「上一轮已生成/失败/澄清，这一轮用户补一句话」时系统应继续报告、转咨询、切专家还是重新澄清。上周已开始统一 router prompt 边界，本周继续把路由规则、工具描述、手测记录和召回脚本补齐，让实际会话更可测。

\textbf{方案说明}：本周把业务边界集中在 router skill 和工具描述里，运行时 prompt 只保留上下文恢复原则；同时增加报告召回、车辆专家召回、车牌容错等手测记录和 fixture，用脚本化样例沉淀真实误判场景。

\textbf{功能说明}：
\begin{itemize}
\item 多轮对话承接更自然：用户上一轮在补充信息、改条件或追问时，系统更容易接上上下文，不会频繁「从头来过」。
\item 失败后的体验更顺：查询未命中时更倾向自动回到澄清流程，而不是直接给出生硬或跑偏的结果。
\item 问法容错更好：线路与车牌的常见口语化写法更容易被识别，减少「明明是同一个对象却查不到」的情况。
\item 报告、专家、咨询的分流更清晰：用户问诊断建议时不容易被直接拉到正式报告，问正式报告时也不容易被分流到泛咨询。
\end{itemize}

\textbf{相关提交（日期）}：
\begin{itemize}
\item \texttt{b9671a1}（2026-04-16）unify router prompt boundaries
\item \texttt{1186be6}（2026-04-17）fix route report partition plumbing
\item \texttt{370b3f5}（2026-04-17）add temporary route lookup normalizer
\item \texttt{bb391c7}（2026-04-18）switch vehicle report input to numberPlate
\item \texttt{2bba37a}（2026-04-21）clarify incomplete Guangdong plates before lookup
\end{itemize}

\section{专家运行时扩展：线路专家接入，专家注册与上下文构建抽象化}

\textbf{背景}：车辆专家已经具备一定召回和咨询能力，但线路类问题仍主要依赖报告或通用咨询承接，容易出现「用户只是想问线路风险，但系统直接出正式报告」的边界问题。本周将线路专家接入专家体系，并同步整理专家上下文构建和注册机制。

\textbf{方案说明}：新增线路专家 skill 与 runtime 支持，把专家能力从单点车辆专家扩展为可注册的专家集合；同时在 chat-service、router-service、worker-runner 和工具校验侧接入 \texttt{route\_expert}，让 router 能把线路咨询、线路风险解释和正式线路报告区分开。

\textbf{功能说明}：
\begin{itemize}
\item 新增了线路专家能力：用户针对线路风险、线路表现的提问，得到的答复更专业、更聚焦。
\item 专家回答边界更清楚：咨询类问题和正式报告请求区分更稳定，减少「答非所问」或越界输出。
\item 专家服务稳定性提升：在真实多轮会话中，专家回答的连续性更好，不容易中途跳转到不相关路径。
\end{itemize}

\textbf{相关提交（日期）}：
\begin{itemize}
\item \texttt{3f9fc7b}（2026-04-20）Add route expert runtime support
\end{itemize}

\section{事故与 ?????? 补齐，同时移除旧 query\_data 报告管线残留}

\textbf{背景}：报告管线切到 MCP 后，如果事故调查、???? gateway 或旧 \texttt{query\_data} 路径仍保留半套逻辑，系统会在不同报告之间出现「有的走 MCP、有的走旧查询、有的 demo 不一致」的状态。这个问题会直接影响本地回归和测试服联调。

\textbf{方案说明}：本周新增事故调查 ?????? 适配，调整 ???? registry/gateway/fixtures，并继续删除旧 \texttt{query\_data} 管线残留。结构上让报告预取、?????? 和 normalizer 尽量走同一类 source 约定。

\textbf{功能说明}：
\begin{itemize}
\item 事故、线路、单位等报告在测试和线上的表现更接近，减少「测试能过、线上不一致」的体验问题。
\item 报告生成链路更统一，用户在不同报告类型间切换时，整体交互风格更稳定。
\item 历史遗留分支持续减少，报告结果的波动和偶发差异进一步下降。
\end{itemize}

\textbf{相关提交（日期）}：
\begin{itemize}
\item \texttt{5ed17d3}（2026-04-17）remove query\_data pipeline and add accident demo mcp
\end{itemize}

\section{文档 chores/算力评估/测试 manual 整理}

\begin{itemize}
\item 扫描版 PDF 处理能力增强：图片型文档可纳入报告与知识处理流程，上传后「无法解析」阻断场景明显减少。
\item 部署与验收路径更清晰：OCR 与报告链路有了可执行的部署、联调和验证说明，交付过程更可控。
\item 问题复现与回归效率提升：关键场景已沉淀为可复用记录，定位问题和复测速度更快。
\item 客户侧算力建议已形成统一口径：当前阶段可继续用 \texttt{4×RTX 3090} 完成测试、联调和验收，不作为正式生产承载。
\item 正式生产建议起点明确：建议至少 \texttt{4×H100} 作为生产入门配置，用于承接在线报告生成和稳定服务能力。
\item 扩容路线明确：当并发、上下文长度和 SLA 目标提升时，再按业务增长评估 \texttt{8×H100} 或 \texttt{H200/B200} 级别方案。
\item 对外沟通结论可直接复用：测试阶段强调「先跑通与验收」，生产阶段强调「稳定承载与可承诺服务」。
\end{itemize}

\textbf{相关提交（日期）}：
\begin{itemize}
\item \texttt{d5084bf}（2026-04-16）show message sources in chat bubble
\item \texttt{8f735db}（2026-04-21）remove dashboard demo and default to assistant
\item \texttt{888ad8a}（2026-04-21）fix unit trend source and suggestion punctuation
\item \texttt{34aa1d1}（2026-04-20）fix unit trend fallback and empty suggestions
\item \texttt{e22cfea}（2026-04-20）fix unit report mcp suggestions and trend
\end{itemize}

\section{待跟进事项}

\begin{itemize}
\item 将 \texttt{scripts/report-source-contract.mjs} 接入固定回归流程（本地脚本或 CI），作为报告来源结构检查。
\item 基于本周新增 fixture 和 manual 记录，继续执行报告召回、专家召回、车牌容错回归。
\item 按 \texttt{OCR\_RAG\_DEPLOYMENT\_MANUAL.md} 在测试环境完成扫描 PDF 的端到端验证并记录结果。
\item 梳理 OCR/RAG 本周新增内容的提交边界，区分源码、文档与依赖/构建产物。
\end{itemize}

\textbf{相关提交（日期）}：本节为待跟进事项，无对应本周单一提交。

\section{下周计划}

\begin{enumerate}
\item 补齐全部管线专家：基于 \texttt{expert-runtime-service-design-20260418.md} 的 5 个 domain 目标（驾驶员/车辆/单位/线路/事故），在当前已具备驾驶员、车辆、线路专家咨询能力的基础上，补齐单位与事故的专家咨询链路，并统一专家与报告之间的边界标准。现状依据是 \texttt{agent/src/domains/experts/registry.ts} 目前仍缺 \texttt{unit + consult}、\texttt{incident + consult}；下周将同步完善对应技能描述、路由分流与多轮承接，避免「该走专家却走报告」或「该走报告却走泛咨询」。验收标准以 4/23 批次 manual 手册与现有召回用例为基线，重点看多轮追问下是否稳定命中正确专家。

\item 迭代 rule response 管线的召回和回复效果质量：围绕「规则命中后是否应走 \texttt{rule\_reply}」和「回复是否保持规则边界」做专项优化，优先处理历史问题中反复出现的「router 代答业务内容」与规则/咨询分流不稳。现状依据包括 \texttt{omni-router-eval-findings-20260316.md}、\texttt{规则列表系统开发文档.md}、\texttt{router-prompt-followups-20260416.md}，以及当前 \texttt{RULE\_MATCH\_RESULTS} / \texttt{RULE\_ROUTING\_POLICY} 注入机制；下周目标是让高置信规则命中、边界分流、\texttt{rule\_exit} 回退后的二次路由更一致。验收方式使用规则场景回归样例与线上真实会话抽样，关注误召回率、错答率和同类问题答复一致性。

\item 上传更多文档测试并迭代 RAG 质量，专家管线接入 RAG：在现有 \texttt{rag-omni-verification-cases-20260412.md} 三类制度文档验证基础上，扩充更多真实文档（含扫描件）做检索与回答质量测试，并按 \texttt{OCR\_RAG\_DEPLOYMENT\_MANUAL.md} 完成 OCR 场景端到端验证，确保扫描 PDF 不再成为入库与问答阻断点。并行推进专家管线接入 RAG 上下文，让专家回答不只依赖结构化数据，也能在制度/流程类问题上引用知识库信息。承载策略继续对齐 \texttt{客户版\_报告生成算力建议说明.md}：测试阶段沿用 \texttt{4×RTX 3090} 做联调与验收，生产能力评估按 \texttt{4×H100} 作为起步标准推进。
\end{enumerate}

\textbf{相关提交（日期）}：本节为下周计划，无对应本周单一提交。

\end{document}
